\chapter{État de l'art}
\label{ch:etat_art}

\section*{Introduction}
\phantomsection
\addcontentsline{toc}{section}{Introduction}

Avant de concevoir quoi que ce soit, il faut
comprendre ce qui existe. Ce chapitre passe en
revue les travaux scientifiques qui touchent
aux quatre dimensions de ce projet : la détection
d'objets en temps réel, le déploiement sur
dispositif embarqué, la protection des données
personnelles, et la décision adaptative des feux
de circulation.

Les articles retenus ont été sélectionnés sur
IEEE Xplore, Google Scholar et Scopus, avec un
focus sur les publications récentes (2024--2026)
indexées dans des journaux Q1 et Q2. Les
références ont été gérées avec Zotero tout au
long du projet.

Ce chapitre s'organise autour de six thèmes :
les smart cities et la gestion du trafic urbain,
la détection d'objets, l'architecture YOLO,
l'Edge AI, la protection des données, et la
décision adaptative. Il se conclut par une
synthèse qui positionne ce travail par rapport
à l'existant et identifie le gap que ce mémoire
tente de combler.

\section{Smart City et gestion du trafic urbain}

\subsection{Définition et concept}

Une ville intelligente exploite les technologies
numériques — capteurs IoT, intelligence
artificielle, analyse de données en temps réel
— pour améliorer la qualité de vie de ses
habitants et résoudre les problèmes urbains
chroniques : congestion, pollution, consommation
énergétique, sécurité~\cite{mystakidis2025}.

\begin{figure}[H]
\centering
\includegraphics[width=0.85\textwidth]
  {smartcity_dimensions}
\caption{Les composantes d'une ville intelligente}
\label{fig:smartcity_dimensions}
\end{figure}

Ce n'est pas un concept abstrait. À Agadir,
la transformation en ville intelligente est
inscrite dans un programme officiel depuis 2022,
avec des axes concrets qui touchent directement
à la mobilité urbaine.

\subsection{Les six dimensions d'une smart city}

La littérature identifie six dimensions
complémentaires qui structurent le concept
de ville intelligente~\cite{hazarika2024}.
La première est la \textbf{Smart Mobility} :
systèmes de transport connectés, régulation
du trafic en temps réel, réduction de la
congestion. La deuxième est le \textbf{Smart
Environment} : gestion durable des ressources,
réduction des émissions. Viennent ensuite la
\textbf{Smart Economy}, la \textbf{Smart
Governance}, le \textbf{Smart Living} et
les \textbf{Smart People} — qui désigne
l'implication active des habitants dans
les transformations urbaines.

Notre travail s'inscrit directement dans
la dimension Smart Mobility, qui conditionne
la fluidité urbaine, la qualité de l'air
et la sécurité des usagers.

\subsection{Le Maroc et l'initiative Smart City
            Agadir}

Le Maroc s'est engagé dans une stratégie
nationale de transformation numérique qui
touche plusieurs secteurs : mobilité,
sécurité urbaine, énergie et services
numériques~\cite{elalaoui2025}. Agadir a
lancé en 2022 son programme officiel de
transformation, structuré autour de douze
axes dans le cadre du Plan d'Action Communal
2022--2027.

\begin{figure}[H]
\centering
\includegraphics[width=0.85\textwidth]
  {agadir_smartcity}
\caption{Axes du Plan d'Action Communal
         d'Agadir 2022--2027
\textit{(Source : Agadir Green City
         Action Plan, Ville d'Agadir, 2022)}}
\label{fig:agadir_plan}
\end{figure}

Une étude menée auprès de 98 jeunes résidents
d'Agadir montre que la majorité soutient cette
transformation, avec une adhésion fortement
corrélée au niveau d'éducation et à la maîtrise
du numérique~\cite{elalaoui2025}. C'est dans
ce contexte favorable que s'inscrit le système
développé dans ce mémoire.

\subsection{La gestion du trafic comme
            composante centrale}

La mobilité est l'une des premières dimensions
visibles d'une ville intelligente, parce qu'elle
affecte directement la vie quotidienne des
habitants. Les feux à cycles fixes — encore
majoritaires dans les villes marocaines —
fonctionnent selon des durées prédéfinies,
sans tenir compte de ce qui se passe réellement
à l'intersection. Résultat : des files d'attente
inutiles quand une voie est vide, et des
blocages quand le trafic dépasse les hypothèses
de dimensionnement du cycle.

\begin{table}[H]
\centering
\caption{Feux à cycles fixes vs système
         adaptatif intelligent}
\begin{tabularx}{\textwidth}{Xcc}
\toprule
\textbf{Critère} &
\textbf{Feux fixes} &
\textbf{Système adaptatif} \\
\midrule
Temps d'attente moyen &
Élevé & Réduit \\
Adaptation au trafic réel &
Non & Oui \\
Priorité véhicules urgence &
Non & Oui \\
Gestion heures de pointe &
Inefficace & Optimisée \\
Détection piétons &
Non & Oui \\
Conformité Loi 09-08 &
Non applicable & Intégrée \\
\bottomrule
\end{tabularx}
\label{tab:feux_comparaison}
\end{table}

Hazarika et al. (2024) confirment qu'un système
basé sur la détection YOLO permet de réduire
significativement les temps d'attente aux
intersections~\cite{hazarika2024}. La
Figure~\ref{fig:traffic_ai} illustre un exemple
de détection et comptage par zones — principe
au cœur de notre approche.

\begin{figure}[H]
\centering
\includegraphics[width=0.90\textwidth]
  {traffic_intersection_ai}
\caption{Détection et comptage des véhicules
         et piétons par zones à une intersection
\textit{(Source : Parquery, 2024)}}
\label{fig:traffic_ai}
\end{figure}

\section{Détection d'objets en temps réel}

\subsection{Définition et enjeux}

La détection d'objets consiste à répondre
simultanément à deux questions sur une image :
\textit{qu'est-ce que c'est ?} et
\textit{où est-ce ?}~\cite{ramos2025}.
C'est différent de la simple classification,
qui identifie sans localiser.

Dans notre cas, on a besoin des deux en même
temps : identifier si un objet est une voiture,
un piéton ou une ambulance, et savoir dans
quelle zone de l'intersection il se trouve
pour alimenter la décision MDP. Et tout ça
en temps réel, sur un CPU ARM sans GPU dédié.

\subsection{Approches one-stage et two-stage}

Il existe deux grandes familles d'algorithmes
de détection. Les approches \textbf{two-stage}
— comme Faster R-CNN — fonctionnent en deux
étapes : d'abord générer des régions candidates,
puis les classifier. Cette double passe donne
une bonne précision, mais au prix d'une latence
incompatible avec le temps réel~\cite{ramos2025}.

\begin{figure}[H]
\centering
\includegraphics[width=0.95\textwidth]{1vs2_stages}
\caption{Comparaison des approches one-stage
         et two-stage
\textit{(Source : Sharkyun, Medium, 2024)}}
\label{fig:detection_approaches}
\end{figure}

Les approches \textbf{one-stage} — YOLO, SSD,
RetinaNet — font tout en un seul passage :
localisation et classification simultanément.
C'est plus rapide, avec une légère concession
sur la précision. C'est la seule option viable
pour un déploiement sur Raspberry Pi 5.

\begin{figure}[H]
\centering
\includegraphics[width=0.92\textwidth]
  {2stagedetection}
\caption{Fonctionnement d'une approche two-stage
\textit{(Source : Sharkyun, Medium, 2024)}}
\label{fig:two_stage_detection}
\end{figure}

\begin{figure}[H]
\centering
\includegraphics[width=0.92\textwidth]
  {1stagedetection}
\caption{Fonctionnement d'une approche one-stage
\textit{(Source : Sharkyun, Medium, 2024)}}
\label{fig:one_stage_detection}
\end{figure}

\begin{table}[H]
\centering
\caption{Comparaison one-stage et two-stage}
\begin{tabularx}{\textwidth}{Xcc}
\toprule
\textbf{Critère} &
\textbf{Two-Stage} &
\textbf{One-Stage} \\
\midrule
Vitesse d'inférence &
Faible & \textbf{Élevée} \\
Précision (mAP) &
\textbf{Élevée} & Bonne \\
Temps réel sur CPU ARM &
Impossible & \textbf{Oui} \\
Déploiement embarqué &
Non & \textbf{Oui} \\
Exemples &
Faster R-CNN & \textbf{YOLO, SSD} \\
\bottomrule
\end{tabularx}
\label{tab:one_two_stage}
\end{table}

\subsection{Métriques d'évaluation}

Plusieurs métriques standardisées sont utilisées
pour évaluer les modèles de détection. On les
retrouve tout au long de ce mémoire, il est
utile de les définir ici.

\subsubsection{Intersection over Union (IoU)}

L'IoU mesure le chevauchement entre la boîte
prédite $\hat{B}$ et la boîte réelle $B$ :

\begin{equation}
\text{IoU} = \frac{|\hat{B} \cap B|}
             {|\hat{B} \cup B|}
\label{eq:iou}
\end{equation}

Une détection est considérée correcte
si $\text{IoU} \geq 0.5$.

\subsubsection{Précision et Rappel}

La \textbf{précision} mesure combien de
détections sont justes parmi toutes celles
produites :

\begin{equation}
\text{Précision} = \frac{VP}{VP + FP}
\label{eq:precision}
\end{equation}

Le \textbf{rappel} mesure combien d'objets
réels ont été détectés :

\begin{equation}
\text{Rappel} = \frac{VP}{VP + FN}
\label{eq:recall}
\end{equation}

Ces deux métriques sont souvent en tension :
augmenter le rappel tend à réduire la
précision, et inversement.

\subsubsection{Mean Average Precision (mAP)}

La mAP est la métrique de référence pour
comparer les modèles de détection. Elle calcule
la moyenne des précisions moyennes sur toutes
les classes $C$ :

\begin{equation}
\text{mAP} = \frac{1}{C} \sum_{c=1}^{C} AP_c
\label{eq:map}
\end{equation}

Dans ce travail, on utilise la \textbf{mAP@50}
(IoU $\geq$ 0.5) et la \textbf{mAP@50-95}
(moyenne sur IoU de 0.5 à 0.95, pas de 0.05).
La mAP@50-95 est plus stricte et pénalise
les localisations imprécises.

\subsubsection{Non-Maximum Suppression (NMS)}

Le NMS est un post-traitement qui élimine
les détections redondantes sur le même objet.
Pour chaque classe, il garde uniquement la
boîte avec la confiance la plus haute parmi
celles qui se chevauchent trop :

\begin{equation}
\hat{B}^* = \arg\max_{\hat{B}}
\text{conf}(\hat{B})\
\text{s.t.}\
\text{IoU}(\hat{B}, \hat{B}') < \theta_{NMS}
\label{eq:nms}
\end{equation}

\begin{definitionbox}{NMS-Free dans YOLO26n}
YOLO26n supprime l'étape NMS grâce à un
mécanisme d'assignation One-to-One introduit
dans YOLOv10. Sur CPU ARM, supprimer ce
post-traitement réduit directement la latence
d'inférence — un avantage concret pour
le déploiement sur Raspberry Pi 5~\cite{sapkota2025}.
\end{definitionbox}

\section{Architecture YOLO : évolution et
         innovations}

\subsection{Principe fondamental}

YOLO (\textit{You Only Look Once}) a été
introduit par Redmon et al. en 2015 avec
une idée simple : au lieu de parcourir
l'image en plusieurs étapes, on la traite
en une seule passe~\cite{ramos2025}.

Le réseau divise l'image en une grille
de $S \times S$ cellules. Pour chaque
cellule, il prédit simultanément les boîtes
englobantes et les probabilités de classe.
La prédiction finale pour chaque boîte est :

\begin{equation}
\text{Score} = P(\text{classe}_i | \text{objet})
\times \text{IoU}_{pred}^{truth}
\label{eq:yolo_score}
\end{equation}

\begin{figure}[H]
\centering
\includegraphics[width=0.88\textwidth]{yoloxyz}
\caption{Détection YOLO : grille et prédiction
         des boîtes englobantes
\textit{(Source : Blent.ai, 2024)}}
\label{fig:yolo_grid}
\end{figure}

\subsection{Architecture générale}

Toutes les versions modernes de YOLO partagent
la même structure en trois blocs :

\begin{figure}[H]
\centering
\includegraphics[width=0.95\textwidth]
  {yoloarchitecture}
\caption{Architecture générale d'un modèle
         YOLO moderne
\textit{(Source : MDPI Electronics, 2025)}}
\label{fig:yolo_architecture}
\end{figure}

Le \textbf{Backbone} extrait les
caractéristiques visuelles de l'image à
différents niveaux d'abstraction. Dans
YOLO26n, il utilise les blocs C3k2 et C2PSA.
Le \textbf{Neck} fusionne ces caractéristiques
à différentes résolutions via FPN et PAN,
ce qui permet de détecter des objets de
tailles variées — petits piétons comme
grands camions. La \textbf{Head} produit
les prédictions finales : coordonnées,
scores de confiance et classes. Les
architectures récentes utilisent une
\textit{Decoupled Head} qui sépare les
tâches de classification et de régression
pour de meilleures performances.

\subsection{Évolution de YOLOv1 à YOLO26}

En dix ans, la famille YOLO a évolué
considérablement, chaque version apportant
des innovations architecturales mesurables
sur des benchmarks standardisés~\cite{ramos2025}.

\begin{table}[H]
\centering
\caption{Évolution des principales versions
         YOLO (2015--2025)}
\begin{tabularx}{\textwidth}{llX}
\toprule
\textbf{Version} &
\textbf{Année} &
\textbf{Innovation principale} \\
\midrule
YOLOv1 & 2015 &
Détection unifiée en une passe,
grille $7\times7$ \\
YOLOv2 & 2016 &
Anchor boxes, entraînement multi-échelle \\
YOLOv3 & 2018 &
Détection à trois niveaux de résolution \\
YOLOv4 & 2020 &
CSP, augmentation Mosaic, CIoU Loss \\
YOLOv5 & 2020 &
Migration PyTorch, variantes
nano/small/medium/large \\
YOLOv7 & 2022 &
E-ELAN, meilleur compromis
vitesse/précision \\
YOLOv8 & 2023 &
Anchor-free, Decoupled Head,
architecture Ultralytics unifiée \\
YOLOv9 & 2024 &
GELAN, Programmable Gradient
Information \\
YOLOv10 & 2024 &
Suppression du NMS, assignation
One-to-One \\
YOLOv11 & 2024 &
C3k2, PSA, amélioration backbone \\
\rowcolor{PrimaryBlue!15}
\textbf{YOLO26} & \textbf{2025} &
\textbf{ProgLoss, STAL, MuSGD,
NMS-Free natif, optimisé edge} \\
\bottomrule
\end{tabularx}
\label{tab:yolo_evolution}
\end{table}

\subsection{Pourquoi YOLO26n a été retenu}

Lors du benchmark, YOLOv8n et YOLOv11n
affichaient une précision parfois proche
ou légèrement meilleure que YOLO26n. Le
choix n'était donc pas évident si on
regardait uniquement les métriques GPU.

C'est le benchmark Raspberry Pi 5 qui a
tranché. Sur CPU ARM, la différence devient
visible :

\begin{table}[H]
\centering
\caption{Comparaison embarquée des variantes
         nano sur Raspberry Pi 5}
\begin{tabularx}{\textwidth}{lcccc}
\toprule
\textbf{Modèle} &
\textbf{Format} &
\textbf{FPS Pi 5} &
\textbf{mAP50} &
\textbf{Params} \\
\midrule
YOLOv8n  & FP32 & 5.3 & 69.8\% & 3.0M \\
YOLOv11n & FP32 & 5.9 & 69.2\% & 2.6M \\
\rowcolor{PrimaryBlue!15}
\textbf{YOLO26n} &
\textbf{FP32} & \textbf{7.6} &
65.6\% & \textbf{2.4M} \\
\bottomrule
\end{tabularx}
\label{tab:yolo_benchmark}
\end{table}

YOLO26n est retenu pour trois raisons
concrètes. Premièrement, il atteint 7.6 FPS
réels sur le Pi 5 — la marge la plus
importante au-dessus du seuil minimal de
5 FPS, ce qui laisse des ressources CPU
disponibles pour les autres modules du
pipeline. Deuxièmement, son architecture
NMS-Free réduit directement la latence sur
CPU ARM, là où les autres modèles paient
encore le coût du post-traitement.
Troisièmement, avec 2.4M de paramètres et
environ 9.8 MB en ONNX FP32, son empreinte
mémoire est la plus faible des trois.

YOLO26n intègre aussi trois innovations
spécifiques qui l'ont rendu attractif pour
ce projet~\cite{sapkota2025} : ProgLoss
pour une convergence plus stable à
l'entraînement, STAL (\textit{Small
Target-Aware Learning}) pour améliorer la
détection des petits objets — piétons et
motos notamment —, et MuSGD, un optimiseur
hybride inspiré des grands modèles de
langage qui accélère la convergence jusqu'à
47\% par rapport aux optimiseurs classiques.

\section{Edge AI et déploiement embarqué}

\subsection{Définition et positionnement}

L'Edge AI consiste à exécuter les modèles
d'intelligence artificielle directement
sur le dispositif qui collecte les données,
sans envoyer quoi que ce soit vers un
serveur distant. Dans notre architecture,
c'est à la fois un choix technique et une
nécessité juridique~\cite{ficili2025}.

\begin{figure}[H]
\centering
\resizebox{0.90\textwidth}{!}{%
\begin{tikzpicture}[
  font=\footnotesize,
  stack/.style={
    draw, rounded corners=7pt,
    align=center, minimum width=4.2cm,
    minimum height=0.88cm, very thick
  },
  cloud/.style={stack,
    draw=red!65, fill=red!8},
  edge/.style={stack,
    draw=AccentGreen, fill=AccentGreen!8},
  title/.style={
    font=\bfseries\small, align=center},
  choice/.style={
    draw=PrimaryBlue, fill=PrimaryBlue!6,
    rounded corners=6pt, align=center,
    text width=9.4cm,
    minimum height=0.95cm,
    font=\scriptsize\bfseries
  },
  note/.style={
    draw=DarkGray!45, fill=white,
    rounded corners=4pt, align=center,
    font=\scriptsize, text width=4.3cm
  },
  arrow/.style={
    -{Latex[length=3mm]},
    line width=1pt}
]

\node[title, color=red!70] at (0,0)
  {Architecture Cloud};
\node[cloud] (c1) at (0,-1.00) {Caméra};
\node[cloud] (c2) at (0,-2.25)
  {Transmission réseau};
\node[cloud] (c3) at (0,-3.50)
  {Serveur distant};
\node[note, draw=red!45] at (0,-4.85)
  {Images brutes transmises hors de
   l'intersection. Risque juridique réel.};

\node[title, color=AccentGreen] at (7.2,0)
  {Architecture Edge AI retenue};
\node[edge] (e1) at (7.2,-1.00) {Caméra};
\node[edge, draw=PrimaryBlue,
      fill=PrimaryBlue!8] (e2) at (7.2,-2.25)
  {Raspberry Pi 5\\YOLO26n FP32};
\node[edge] (e3) at (7.2,-3.50)
  {Décision locale};
\node[note, draw=AccentGreen!55] at (7.2,-4.85)
  {Traitement local, aucune donnée
   personnelle ne quitte l'intersection.};

\draw[arrow, red!65] (c1) -- (c2);
\draw[arrow, red!65] (c2) -- (c3);
\draw[arrow, AccentGreen] (e1) -- (e2);
\draw[arrow, AccentGreen] (e2) -- (e3);

\node[choice] at (3.6,-6.25)
  {Choix retenu : traitement local,
   latence faible, conformité Loi 09-08.};

\end{tikzpicture}}
\caption{Architecture Cloud vs Edge AI
         pour la gestion du trafic}
\label{fig:cloud_vs_edge}
\end{figure}

Un système cloud transmet les flux vidéo
vers des serveurs distants — ce qui revient
à envoyer en continu des images contenant
des visages et des plaques d'immatriculation
hors de l'intersection. C'est exactement
ce que la Loi 09-08 encadre strictement.
Le traitement local sur le Pi 5 règle ce
problème à la racine.

\subsection{Défis spécifiques du Raspberry Pi 5}

Déployer un pipeline de vision par ordinateur
sur CPU ARM soulève des contraintes absentes
des environnements GPU d'entraînement~\cite{wang2025}.

Le processeur ARM Cortex-A76 n'a pas d'unité
de calcul tensoriel : toute l'inférence repose
sur le CPU. La RAM de 8 Go est partagée entre
le système, le modèle et le pipeline vidéo.
La dissipation thermique est limitée — une
charge CPU soutenue peut déclencher le
\textit{thermal throttling} et réduire les
performances. Et enfin, la gestion du trafic
exige au minimum 5 FPS pour que les décisions
restent cohérentes avec l'évolution réelle
de l'intersection.

\subsection{Optimisations retenues}

\subsubsection{Variante nano uniquement}

Le benchmark a confirmé empiriquement que
les variantes \textit{small} ne passent pas
la barre des 3 FPS sur Raspberry Pi 5, quel
que soit le format d'export. Les variantes
nano sont les seules compatibles avec le
temps réel sur CPU ARM~\cite{wang2025}.

\subsubsection{Format FP32 retenu, INT8 écarté}

La quantification INT8 réduit la taille du
modèle, mais elle n'apporte pas de gain FPS
systématique sur Raspberry Pi 5 et introduit
une confusion plus marquée entre les classes
visuellement proches — notamment car, bus et
truck. Le format FP16 est incompatible avec
le CPU ARM Cortex-A76 : sans Tensor Cores,
les opérations FP16 sont émulées en logiciel
et ralentissent l'inférence. Le format FP32
est donc retenu.

\begin{table}[H]
\centering
\caption{Formats de précision sur
         Raspberry Pi 5 — modèle de base}
\begin{tabularx}{\textwidth}{lccXl}
\toprule
\textbf{Format} &
\textbf{FPS} &
\textbf{Latence} &
\textbf{Observation} &
\textbf{Décision} \\
\midrule
FP32 & 3.7 & 178ms &
Référence stable &
Retenu \\
FP16 & — & — &
Incompatible CPU ARM &
Écarté \\
INT8 & 5.0 & 193ms &
Confusion car/bus/truck &
Écarté \\
\bottomrule
\end{tabularx}
\label{tab:formats}
\end{table}

\subsubsection{Export ONNX}

Le modèle PyTorch natif est optimisé pour
GPU et peu efficace sur CPU ARM. L'export
au format ONNX (\textit{Open Neural Network
Exchange}) permet d'utiliser ONNX Runtime
avec le fournisseur \texttt{CPUExecutionProvider},
conçu pour maximiser l'inférence sur
processeur~\cite{wang2025} :

\begin{equation}
\texttt{best.pt}
\xrightarrow{\text{export}}
\texttt{yolo26n.onnx}
\xrightarrow{\text{ONNX Runtime}}
\text{Inférence Pi 5}
\label{eq:onnx}
\end{equation}

\section{Protection des données personnelles}

\subsection{Cadre légal marocain — Loi 09-08}

La Loi n°09-08 du 18 février 2009 est le texte
de référence au Maroc pour la protection des
données personnelles. Elle a institué la
Commission Nationale de contrôle de la
protection des Données à caractère Personnel
(CNDP), dont la mission est de vérifier que
tout traitement de données est licite et ne
porte pas atteinte aux droits fondamentaux
des citoyens~\cite{elalaoui2025}. Le texte
complet est disponible sur le site officiel
de la CNDP.

\begin{figure}[H]
\centering
\includegraphics[width=0.55\textwidth]{cndp}
\caption{Logo de la CNDP, autorité marocaine
         de protection des données
\textit{(Source : cndp.ma, 2024)}}
\label{fig:cndp}
\end{figure}

La loi repose sur quatre principes fondamentaux.
La \textbf{finalité} : les données ne peuvent
être utilisées qu'aux fins déclarées.
La \textbf{proportionnalité} : on ne collecte
que ce qui est strictement nécessaire.
La \textbf{sécurité} : des mesures techniques
adaptées doivent protéger les données.
Et la \textbf{durée limitée} : on ne conserve
pas les données au-delà du nécessaire.

\subsection{Enjeux dans les systèmes de
            surveillance du trafic}

Une caméra installée à un carrefour ne filme
pas que des voitures. Elle capture aussi des
visages — donnée biométrique sensible —, des
plaques d'immatriculation qui permettent
d'identifier le propriétaire du véhicule,
et des patterns de déplacement qui révèlent
les habitudes de vie. Ce sont toutes des
données personnelles au sens de la Loi 09-08.

Des travaux récents montrent que même des
images partiellement anonymisées restent
vulnérables à des attaques de ré-identification
par des algorithmes d'IA avancés~\cite{razi2025}.
Flouter a posteriori ne suffit pas si les
données brutes ont déjà circulé dans le
système.

\begin{figure}[H]
\centering
\resizebox{0.92\textwidth}{!}{%
\begin{tikzpicture}[
  font=\small,
  cell/.style={
    draw, rounded corners=4pt,
    align=center, minimum width=4.05cm,
    minimum height=0.72cm,
    font=\scriptsize
  },
  risk/.style={cell,
    draw=red!60, fill=red!8},
  sol/.style={cell,
    draw=AccentGreen, fill=AccentGreen!8},
  head/.style={
    draw, rounded corners=5pt,
    align=center, minimum width=4.05cm,
    minimum height=0.72cm,
    font=\small\bfseries, text=white
  },
  arrow/.style={
    -{Latex[length=2.7mm]},
    line width=0.95pt, PrimaryBlue}
]

\node[head, fill=red!68] at (0,0) {Risque};
\node[head, fill=PrimaryBlue] at (5.0,0)
  {Mesure technique};
\node[head, fill=AccentGreen!80!black]
  at (10.0,0) {Effet obtenu};

\node[risk] (r1) at (0,-0.95)
  {Identification directe};
\node[sol]  (m1) at (5.0,-0.95)
  {Floutage gaussien\\des personnes};
\node[sol]  (e1) at (10.0,-0.95)
  {Visages non exploitables};

\node[risk] (r2) at (0,-1.85)
  {Suivi des déplacements};
\node[sol]  (m2) at (5.0,-1.85)
  {Tracking après anonymisation};
\node[sol]  (e2) at (10.0,-1.85)
  {Identité non conservée};

\node[risk] (r3) at (0,-2.75)
  {Fuite de données};
\node[sol]  (m3) at (5.0,-2.75)
  {Traitement local Edge AI};
\node[sol]  (e3) at (10.0,-2.75)
  {Aucun envoi Cloud};

\node[risk] (r4) at (0,-3.65)
  {Non-conformité Loi 09-08};
\node[sol]  (m4) at (5.0,-3.65)
  {Anonymisation by design};
\node[sol]  (e4) at (10.0,-3.65)
  {Privacy-by-Design};

\foreach \r/\m/\e in
  {r1/m1/e1,r2/m2/e2,r3/m3/e3,r4/m4/e4}{
  \draw[arrow] (\r.east) -- (\m.west);
  \draw[arrow] (\m.east) -- (\e.west);
}

\node[
  draw=AccentOrange, fill=AccentOrange!7,
  rounded corners=5pt, align=center,
  font=\scriptsize, text width=9.6cm
] at (5.0,-4.65)
{Principe retenu : anonymiser avant toute
visualisation ou transmission, et limiter
le système aux statistiques nécessaires
à la décision de trafic.};
\end{tikzpicture}}
\caption{Risques de confidentialité et
         mesures adoptées dans notre système}
\label{fig:privacy_risks}
\end{figure}

\subsection{Techniques d'anonymisation}

Jeremiah et al. (2024) recensent plusieurs
techniques applicables aux flux vidéo de
surveillance~\cite{jeremiah2024}. Le floutage
gaussien est le plus rapide et le plus
compatible avec un déploiement edge. La
pixelisation est rapide aussi, mais reste
réversible par des algorithmes de
super-résolution. Les approches par
silhouette perdent trop d'information visuelle.
Les méthodes GAN sont trop lourdes pour un
Pi 5. La confidentialité différentielle est
complexe à calibrer en production.

\begin{table}[H]
\centering
\caption{Comparaison des techniques
         d'anonymisation vidéo}
\begin{tabularx}{\textwidth}{lllXl}
\toprule
\textbf{Technique} &
\textbf{Vitesse} &
\textbf{Edge} &
\textbf{Limite principale} &
\textbf{Retenu} \\
\midrule
Floutage gaussien &
Rapide & Oui &
Partiellement réversible &
\textbf{Oui} \\
Pixelisation &
Rapide & Oui &
Réversible par SR &
Non \\
Silhouette &
Moyen & Partiel &
Perte d'information &
Non \\
GAN anonymisation &
Lent & Non &
Trop lourd pour Pi 5 &
Non \\
Differential Privacy &
Variable & Partiel &
Complexe à calibrer &
Non \\
\bottomrule
\end{tabularx}
\label{tab:anon_techniques}
\end{table}

\subsection{Privacy-by-Design}

Le Privacy-by-Design est un principe
architectural qui dit que la protection des
données doit être intégrée dès la conception
du système, pas ajoutée en fin de pipeline.
Dans notre architecture, cela se traduit par
trois règles concrètes : le floutage est
appliqué immédiatement après la détection,
avant tout autre traitement ; aucune vidéo
n'est stockée ni transmise ; et l'intégralité
du pipeline s'exécute en local sur le Pi 5.

La limite de cette approche est connue et
documentée : le floutage gaussien reste
théoriquement vulnérable à des techniques
de super-résolution avancées~\cite{razi2025,
jeremiah2024}. Dans notre contexte — où
l'objectif est le comptage et non
l'identification — cette limite est
acceptable. Elle est néanmoins documentée
dans les limites du Chapitre 6.

\section{Décision adaptative des feux de
         circulation}

\subsection{Limites des systèmes à cycles fixes}

Un feu à cycle fixe ignore ce qui se passe
réellement à l'intersection. Il donne le vert
à une direction pendant 30 secondes même si
cette direction est vide, pendant qu'une
ambulance attend de l'autre côté. C'est le
problème concret que ce module cherche à
résoudre.

Li et al. (2025) quantifient les gains
possibles : une approche adaptative réduit
le temps d'attente moyen de 27,34\% par
rapport aux cycles fixes, avec des gains
atteignant 68,40\% dans les cas de trafic
fortement déséquilibré~\cite{li2025}.

\subsection{Formalisation MDP}

Un Processus de Décision Markovien (MDP)
est un cadre mathématique pour modéliser
la prise de décision séquentielle dans un
environnement dont l'état change en réponse
aux actions prises. Il est défini par le
quadruplet $(S, A, T, R)$~\cite{michailidis2025} :
l'espace d'états $S$, l'espace d'actions $A$,
la fonction de transition $T(s,a,s')$ et
la fonction de récompense $R(s,a)$.

\subsubsection{État $S$}

L'état de l'intersection est représenté par
un vecteur à six composantes :

\begin{equation}
s = (q_N,\ q_S,\ q_E,\ q_O,\ p,\ e)
\label{eq:state}
\end{equation}

où $q_N, q_S, q_E, q_O$ sont les densités
de véhicules dans les quatre directions,
$p$ est le nombre de piétons dans la zone
de passage, et $e$ indique la présence
d'un véhicule d'urgence ($e \in \{0,1\}$).

\subsubsection{Actions $A$}

Cinq phases de feux sont disponibles,
chacune associée à une durée variable
selon la densité détectée :

\begin{table}[H]
\centering
\caption{Phases et durées du module de
         décision}
\begin{tabularx}{\textwidth}{llXl}
\toprule
\textbf{Phase} &
\textbf{Feux verts} &
\textbf{Condition} &
\textbf{Durée} \\
\midrule
North-South & Nord + Sud &
Densité NS maximale & 15--45s \\
East-West   & Est + Ouest &
Densité EO maximale & 15--45s \\
Pedestrian  & Piétons &
Piétons + trafic faible & 30s \\
Emergency   & Tous rouges &
Véhicule urgence détecté & 45s \\
All-Red     & Aucun &
Aucun trafic & 15s \\
\bottomrule
\end{tabularx}
\label{tab:phases}
\end{table}

\subsubsection{Récompense $R$}

La récompense est l'opposé du temps
d'attente moyen (AWT) à l'intersection :

\begin{equation}
R(s, a) = -\text{AWT}(s, a)
\label{eq:reward}
\end{equation}

Maximiser la récompense revient à minimiser
les temps d'attente — c'est exactement
l'objectif opérationnel du système.

\subsection{Politique de décision déterministe}

Les approches par Reinforcement Learning
profond donnent de bons résultats en
simulation, mais elles nécessitent l'inférence
d'un réseau de politique en temps réel —
incompatible avec les ressources du
Raspberry Pi 5. On a donc choisi une politique
déterministe basée sur une fonction de
scoring pondéré, calculable en quelques
microsecondes.

Le score pour chaque direction $i$ est :

\begin{equation}
\text{Score}(i) = w_v \cdot q_i^{veh}
               + w_p \cdot q_i^{ped}
               + w_e \cdot q_i^{urg}
\label{eq:score}
\end{equation}

Les coefficients ont été fixés par conception,
selon une logique de priorité trafic, puis
validés sur des scénarios représentatifs.
$w_v = 1.0$ est le poids de base pour les
véhicules. $w_p = 1.5$ donne une priorité
aux piétons sans bloquer tout le trafic
véhicule. $w_e = 100.0$ représente la
priorité absolue aux véhicules d'urgence —
une valeur suffisamment haute pour qu'aucun
autre score ne puisse la dépasser.

\begin{table}[H]
\centering
\caption{Coefficients de pondération MDP}
\begin{tabularx}{\textwidth}{llXl}
\toprule
\textbf{Coeff.} &
\textbf{Valeur} &
\textbf{Justification} &
\textbf{Classes} \\
\midrule
$w_v$ & 1.0 &
Référence véhicule &
car, bus, truck, moto \\
$w_p$ & 1.5 &
Priorité piétons (+50\%) &
person \\
$w_e$ & 100.0 &
Priorité absolue urgence &
emergency\_vehicle \\
\bottomrule
\end{tabularx}
\label{tab:coefficients}
\end{table}

La politique complète suit une hiérarchie
de règles :

\begin{equation}
\pi^*(s) =
\begin{cases}
\text{Emergency} &
\text{si } e > 0 \\[6pt]
\text{Pedestrian} &
\text{si } p > 0 \text{ et }
\max_i q_i^{veh} < \theta_H \\[6pt]
\arg\max_i \text{Score}(i) &
\text{sinon}
\end{cases}
\label{eq:policy}
\end{equation}

La durée de phase dépend de la densité
détectée :

\begin{equation}
D^* =
\begin{cases}
45\text{s} & q_i^{veh} \geq \theta_H = 10 \\
30\text{s} & \theta_M \leq q_i^{veh}
             < \theta_H \\
15\text{s} & q_i^{veh} < \theta_M = 5
\end{cases}
\label{eq:duration}
\end{equation}

\begin{figure}[H]
\centering
\resizebox{0.88\textwidth}{!}{%
\begin{tikzpicture}[
  font=\footnotesize,
  mdpstep/.style={
    draw=PrimaryBlue, fill=PrimaryBlue!7,
    rounded corners=6pt, align=center,
    text width=5.45cm,
    minimum height=1.05cm
  },
  test/.style={
    draw=AccentOrange, fill=AccentOrange!9,
    rounded corners=6pt, align=center,
    text width=5.45cm,
    minimum height=1.05cm
  },
  action/.style={
    draw, rounded corners=6pt,
    align=center, text width=4.15cm,
    minimum height=1.02cm,
    font=\small\bfseries
  },
  arrow/.style={
    -{Latex[length=3mm]},
    line width=1.05pt, color=DarkGray},
  yeslabel/.style={
    font=\scriptsize\bfseries,
    fill=white, inner sep=2pt},
  nolabel/.style={
    font=\scriptsize\bfseries,
    fill=white, inner sep=2pt,
    text=DarkGray}
]

\node[mdpstep] (state) at (0,0)
  {État : $s=(q_N,q_S,q_E,q_O,p,e)$};
\node[test] (urg) at (0,-1.65)
  {Urgence détectée ? $e>0$};
\node[test] (ped) at (0,-3.30)
  {Piétons + trafic faible ?};
\node[mdpstep] (score) at (0,-4.95)
  {Score$(i)=w_vq_i+w_pp_i+w_ee_i$};
\node[test] (traffic) at (0,-6.60)
  {Trafic présent ?};

\node[action, draw=red!65, fill=red!8]
  (a1) at (7.45,-1.65) {Emergency 45s};
\node[action, draw=AccentOrange,
  fill=AccentOrange!10]
  (a2) at (7.45,-3.30) {Pedestrian 30s};
\node[action, draw=AccentGreen,
  fill=AccentGreen!9]
  (a3) at (7.45,-5.72)
  {Phase optimale\\$\arg\max_i$ Score};
\node[action, draw=DarkGray, fill=LightGray]
  (a4) at (7.45,-7.38) {All-Red 15s};

\draw[arrow] (state.south) -- (urg.north);
\draw[arrow] (urg.south) -- (ped.north);
\draw[arrow] (ped.south) -- (score.north);
\draw[arrow] (score.south) -- (traffic.north);
\draw[arrow, red!65] (urg.east) --
  node[yeslabel, above] {Oui} (a1.west);
\draw[arrow, AccentOrange] (ped.east) --
  node[yeslabel, above] {Oui} (a2.west);
\draw[arrow, AccentGreen] (traffic.east)
  -- ++(0.95,0)
  |- node[yeslabel, pos=0.22, above]
     {Oui} (a3.west);
\draw[arrow] (traffic.east) -- ++(1.25,0)
  |- node[nolabel, pos=0.22, below]
     {Non} (a4.west);

\node[
  draw=DarkGray!45, fill=white,
  rounded corners=5pt, align=center,
  font=\scriptsize, text width=11.8cm
] at (3.75,-8.85)
{$w_v=1.0$,\ $w_p=1.5$,\ $w_e=100.0$.
Politique déterministe,
calculable en temps réel sur Pi 5.};
\end{tikzpicture}}
\caption{Organigramme de la politique
         $\pi^*(s)$ — Module M5
         \textit{(Source : auteur)}}
\label{fig:mdp_flowchart}
\end{figure}

\subsection{Validation sur scénarios}

La politique a été testée sur trois scénarios
représentatifs des situations courantes dans
les intersections urbaines. Dans chaque cas,
la phase produite correspond à la phase
attendue par la logique trafic :

\begin{table}[H]
\centering
\caption{Validation de $\pi^*(s)$ sur
         5 scénarios}
\begin{tabularx}{\textwidth}{lXll}
\toprule
\textbf{Scénario} &
\textbf{État détecté} &
\textbf{Phase} &
\textbf{Durée} \\
\midrule
Dense Nord &
$q_N=12$, $p=2$ &
North-South & 45s \\
Trafic équilibré &
4 veh par direction &
North-South & 30s \\
Urgence Est &
$e=1$ &
Emergency & 45s \\
Piétons seuls &
$p=4$, $q<5$ &
Pedestrian & 30s \\
Aucun trafic &
$q=0$, $p=0$ &
All-Red & 15s \\
\bottomrule
\end{tabularx}
\label{tab:validation_mdp}
\end{table}

\section{Synthèse et positionnement}

\subsection{Ce que la littérature fait
            et ce qui manque}

Le tableau suivant positionne les douze
travaux analysés par rapport aux quatre
dimensions de ce projet.

\begin{table}[H]
\centering
\footnotesize
\renewcommand{\arraystretch}{1.12}
\caption{Synthèse comparative de l'état
         de l'art}
\label{tab:synthese}
\begin{tabularx}{\textwidth}
  {p{2.4cm}p{1.15cm}XXXX}
\toprule
\textbf{Travail} &
\textbf{Année} &
\textbf{Détection} &
\textbf{Edge AI} &
\textbf{Privacy} &
\textbf{Décision} \\
\midrule
Mystakidis & 2025 &
Review & — & — & — \\
Hazarika & 2024 &
tiny-YOLOv3 & Raspberry Pi &
Non & SJF+RR \\
El Alaoui & 2025 &
— & — & — &
Étude sociale \\
Ramos & 2025 &
Review YOLO & Partiel & — & — \\
Dinesh & 2024 &
YOLOv7 & PC standard &
Non & Densité \\
Sapkota & 2025 &
YOLO26n & Jetson Orin &
Non & — \\
Wang & 2025 &
Review léger & Pi 4B & — & — \\
Ficili & 2025 &
— & Pi 3 + IoT &
Partiel & Monitoring \\
Jeremiah & 2024 &
MobileNet & Edge MEC &
Floutage & — \\
Razi & 2025 &
— & — & Survey & — \\
Li & 2025 &
Capteurs & — & FL &
Federated DRL \\
Michailidis & 2025 &
Review RL & — & — &
RL Survey \\
\midrule
\rowcolor{PrimaryBlue!15}
\textbf{Notre travail} &
\textbf{2026} &
\textbf{YOLO26n} &
\textbf{Pi 5 8GB} &
\textbf{PbD} &
\textbf{MDP} \\
\bottomrule
\end{tabularx}
\end{table}

Sapkota et al. (2025) utilisent YOLO26n
mais sur Jetson Orin — un hardware beaucoup
plus puissant. Hazarika et al. (2024)
déploient sur Raspberry Pi mais avec
tiny-YOLOv3 et sans protection des données.
Aucun travail ne combine les quatre
dimensions sur un seul dispositif.

\begin{figure}[H]
\centering
\resizebox{0.92\textwidth}{!}{%
\begin{tikzpicture}[
  font=\small,
  pillar/.style={
    draw, rounded corners=7pt,
    text width=3.15cm,
    minimum height=1.22cm,
    align=center,
    font=\small\bfseries,
    very thick
  },
  center/.style={
    draw=DarkGray, fill=white,
    rounded corners=8pt,
    text width=5.6cm,
    minimum height=1.42cm,
    align=center,
    font=\small\bfseries,
    very thick, drop shadow
  },
  gap/.style={
    draw=AccentOrange,
    fill=AccentOrange!8,
    rounded corners=5pt,
    align=center, font=\scriptsize,
    text width=12.2cm
  },
  plus/.style={
    font=\Large\bfseries,
    color=DarkGray},
  arrow/.style={
    -{Latex[length=3mm]},
    line width=1.15pt, color=DarkGray}
]

\node[pillar, draw=PrimaryBlue,
  fill=PrimaryBlue!9]
  (detect) at (0,2.05)
  {Détection\\temps réel\\YOLO26n};
\node[pillar, draw=AccentGreen,
  fill=AccentGreen!9]
  (edge) at (4.05,2.05)
  {Edge AI\\Raspberry Pi 5};
\node[pillar, draw=AccentOrange,
  fill=AccentOrange!10]
  (privacy) at (8.10,2.05)
  {Protection données\\Loi 09-08};
\node[pillar, draw=red!65, fill=red!8]
  (mdp) at (12.15,2.05)
  {Décision\\adaptative\\MDP};

\node[plus] at (2.03,2.05) {$+$};
\node[plus] at (6.08,2.05) {$+$};
\node[plus] at (10.13,2.05) {$+$};

\node[
  draw=DarkGray!55, fill=LightGray!35,
  rounded corners=5pt, align=center,
  font=\scriptsize\bfseries,
  text width=11.8cm
] (combo) at (6.08,0.45)
  {Gap : aucun travail ne combine
  les quatre dimensions sur
  un seul dispositif embarqué.};

\node[center] (work) at (6.08,-1.35)
  {Pipeline M1--M6\\
  \textit{Notre contribution}};

\draw[arrow] (combo.south) -- (work.north);

\node[gap] at (6.08,-3.10)
  {Ce travail combine détection YOLO26n,
  anonymisation Privacy-by-Design,
  décision MDP et déploiement sur Pi 5
  dans un seul pipeline validé sur des
  données incluant le contexte marocain.};

\end{tikzpicture}}
\caption{Positionnement — intersection
         des quatre dimensions}
\label{fig:positioning}
\end{figure}

\begin{table}[H]
\centering
\caption{Différenciation par rapport
         aux travaux les plus proches}
\begin{tabularx}{\textwidth}{lXX}
\toprule
\textbf{Critère} &
\textbf{Travaux existants} &
\textbf{Ce travail} \\
\midrule
Modèle &
YOLOv3 à YOLOv11 &
YOLO26n (2025) \\
Hardware &
PC, Jetson, Pi 3-4 &
Raspberry Pi 5 8GB \\
Format &
PyTorch, TFLite &
ONNX FP32 \\
Privacy &
Absente ou partielle &
Privacy-by-Design \\
Décision &
RL complexe ou fixe &
MDP déterministe \\
Contexte &
Europe, Asie, USA &
Maroc (Agadir) \\
Pipeline &
Partiel &
Complet M1--M6 \\
\bottomrule
\end{tabularx}
\label{tab:positioning}
\end{table}

\section*{Conclusion}
\addcontentsline{toc}{section}{Conclusion}

Ce chapitre a passé en revue les travaux
existants sur les quatre dimensions du
projet : détection, déploiement embarqué,
protection des données et décision
adaptative.

Le constat principal est que chaque dimension
est traitée séparément dans la littérature.
Personne ne les combine dans un pipeline
unique déployé sur dispositif minimal et
validé dans un contexte urbain marocain.
C'est ce gap que le système présenté dans
les chapitres suivants tente de combler,
avec les contraintes et les limites que
cela implique.